% --- Archivo: 04_cono_tangente_primal.tex ---

% =========================================================
\section{Condiciones de optimalidad en forma primal para problemas con restricciones. El cono tangente.}
% =========================================================

Consideramos ahora un problema en formato general
\begin{equation}
    \min f(x) \quad \text{sujeto a} \quad x \in D, \label{eq:problema_general_restricciones}
\end{equation}
donde $D \subset \Rn$ es un conjunto dado (no vacío) y $f : \Rn \to \R$ es la función objetivo. 
Para desarrollar las condiciones de optimalidad necesitamos estudiar la estructura del conjunto factible mediante \textbf{direcciones viables} y \textbf{direcciones tangentes}.

\begin{definition}[Dirección viable]\label{def:direccion_viable}
    Decimos que $d \in \Rn$ es una \textbf{dirección viable} en relación al conjunto $D$ en el punto $\bar{x} \in D$, cuando existe $\varepsilon > 0$ tal que
    \[
        \bar{x} + td \in D \quad \forall t \in [0, \varepsilon].
    \]
    Denotamos por $\mathcal{V}_D(\bar{x})$ el conjunto de todas las direcciones viables en relación al conjunto $D$ en el punto $\bar{x} \in D$.
\end{definition}

\begin{center}
\begin{tikzpicture}[>=Stealth, scale=1.2]
    \draw[fill=MidnightBlue!10, draw=MidnightBlue!80!black, thick] (0,0) -- (3,0) -- (3,2) -- (0.5,2) -- (0,1) -- cycle;
    \filldraw (1,1) circle (1.5pt) node[below] {$\bar{x}$};
    \draw[->, thick, ForestGreen] (1,1) -- (2.5,1.5) node[right] {$d^1$};
    \draw[->, thick, ForestGreen] (1,1) -- (1.5,1.8) node[above] {$d^2$};
    \draw[->, thick, red!80!black] (1,1) -- (0,2) node[left] {$d^3$ (no viable)};
    \node[MidnightBlue!80!black] at (2.5,0.5) {$D$};
\end{tikzpicture}
\end{center}

\begin{definition}[Cono]\label{def:cono}
    Un conjunto $K \subset \Rn$ se llama \textbf{cono} cuando
    \[
        d \in K \implies td \in K \quad \forall t \in \R_+.
    \]
\end{definition}

Es fácil ver que para $D \subset \Rn$ y $\bar{x} \in D$ cualesquiera, el conjunto $\mathcal{V}_D(\bar{x})$ siempre es un cono no vacío (al menos, $0 \in \mathcal{V}_D(\bar{x})$).

\begin{definition}[Dirección de descenso]\label{def:direccion_descenso}
    Decimos que $d \in \Rn$ es una \textbf{dirección de descenso} de $f : \Rn \to \R$ en el punto $\bar{x} \in \Rn$, si existe $\varepsilon > 0$ tal que
    \[
        f(\bar{x} + td) < f(\bar{x}) \quad \forall t \in (0, \varepsilon].
    \]
    Denotamos por $\mathcal{D}_f(\bar{x})$ el conjunto de todas las direcciones de descenso de la función $f$ en el punto $\bar{x}$.
\end{definition}

El conjunto $\mathcal{D}_f(\bar{x})$ puede ser vacío, pero el conjunto $\mathcal{D}_f(\bar{x}) \cup \{0\}$ es siempre un cono.

\begin{lemma}[Direcciones de descenso]\label{lem:caracterizacion_descenso}
    Sea $f : \Rn \to \R$ una función diferenciable en el punto $\bar{x} \in \Rn$. Entonces:
    \begin{enuthm}
        \item Para todo $d \in \mathcal{D}_f(\bar{x})$, se tiene que $\interno{\nabla f(\bar{x})}{d} \le 0$.
        \item Si $d \in \Rn$ satisface $\interno{\nabla f(\bar{x})}{d} < 0$, entonces $d \in \mathcal{D}_f(\bar{x})$.
    \end{enuthm}
\end{lemma}

\begin{proof}
    Sea $d \in \mathcal{D}_f(\bar{x})$. Para todo $t > 0$ suficientemente pequeño, $0 > f(\bar{x} + td) - f(\bar{x}) = t(\interno{\nabla f(\bar{x})}{d} + o(t)/t)$. Pasando al límite cuando $t \to 0^+$, obtenemos $0 \ge \interno{\nabla f(\bar{x})}{d}$, lo que prueba el ítem (a).
    
    Si $\interno{\nabla f(\bar{x})}{d} < 0$, tenemos que $f(\bar{x} + td) - f(\bar{x}) = t(\interno{\nabla f(\bar{x})}{d} + o(t)/t) \le \frac{t}{2}\interno{\nabla f(\bar{x})}{d} < 0$ para todo $t > 0$ suficientemente pequeño, implicando que $d \in \mathcal{D}_f(\bar{x})$.
\end{proof}

Si $\bar{x} \in D$ es un minimizador local de \eqref{eq:problema_general_restricciones}, obviamente $\mathcal{D}_f(\bar{x}) \cap \mathcal{V}_D(\bar{x}) = \emptyset$. Por el \cref{lem:caracterizacion_descenso}, tenemos entonces que $\interno{\nabla f(\bar{x})}{d} \ge 0$ para todo $d \in \mathcal{V}_D(\bar{x})$. Sin embargo, esta condición no siempre proporciona información útil, porque $\mathcal{V}_D(\bar{x})$ puede contener solo al vector nulo. Necesitamos considerar las \textbf{direcciones tangentes}.

\begin{definition}[Dirección tangente y cono tangente]\label{def:cono_tangente}
    Decimos que $d \in \Rn$ es una \textbf{dirección tangente} en relación al conjunto $D$ en el punto $\bar{x} \in D$ cuando
    \[
        \dist(\bar{x} + td, D) = o(t), \quad t \in \R_+.
    \]
    Denotamos por $\mathcal{T}_D(\bar{x})$ el conjunto de todas las direcciones tangentes, el cual forma un cono.
\end{definition}

Claramente, $\mathcal{V}_D(\bar{x}) \subset \mathcal{T}_D(\bar{x})$. Otra noción útil es la del \textbf{cono de Bouligand}:
\begin{equation}
    \mathcal{B}_D(\bar{x}) = \left\{ d \in \Rn \;\middle|\; \begin{aligned} &\exists \{t_k\} \subset \R_+, \, \{t_k\} \to 0^+, \\ &\dist(\bar{x} + t_k d, D) = o(t_k) \end{aligned} \right\}. \label{eq:cono_bouligand}
\end{equation}
De forma equivalente, $d \in \mathcal{B}_D(\bar{x})$ si existen sucesiones $\{t_k\} \to 0^+$ y $\{d^k\} \to d$ tales que $\bar{x} + t_k d^k \in D$ para todo $k$. En general, $\mathcal{T}_D(\bar{x}) \subset \mathcal{B}_D(\bar{x})$. 

\begin{example}
    Sean $q \in (0,1)$ y $D = \{0\} \cup \{q^0, q^1, q^2, \dots\} \subset \R$. En $\bar{x} = 0 \in D$, para todo $d > 0$, eligiendo $t_k = q^k / d$ y $d^k = d$, tenemos $\bar{x} + t_k d^k = q^k \in D$. Así que $d \in \mathcal{B}_D(\bar{x})$ y $\mathcal{B}_D(\bar{x}) = \R_+$. Sin embargo, se puede demostrar analíticamente que $\mathcal{T}_D(\bar{x}) = \{0\}$, probando que $\mathcal{T}_D(\bar{x}) \neq \mathcal{B}_D(\bar{x})$.
\end{example}

\begin{theorem}[Condición necesaria en forma primal]\label{thm:cn_primal}
    Sean $D \subset \Rn$ y $f : \Rn \to \R$ diferenciable en el punto $\bar{x} \in D$. Si $\bar{x}$ es una solución local del problema de minimizar $f$ en $D$, entonces
    \begin{equation}
        \interno{\nabla f(\bar{x})}{d} \ge 0 \quad \forall d \in \mathcal{B}_D(\bar{x}). \label{eq:condicion_primal_primer_orden}
    \end{equation}
\end{theorem}

\begin{proof}
    Fijemos $d \in \mathcal{B}_D(\bar{x}) \setminus \{0\}$ y las sucesiones asociadas $\{t_k\} \to 0^+$ y $\{d^k\} \to d$ tales que $\bar{x} + t_k d^k \in D$. Como $\bar{x}$ es un minimizador local, para todo $k$ suficientemente grande:
    \[
        0 \le f(\bar{x} + t_k d^k) - f(\bar{x}) = t_k \interno{\nabla f(\bar{x})}{d^k} + o(t_k).
    \]
    Dividiendo por $t_k > 0$ y pasando al límite, obtenemos \eqref{eq:condicion_primal_primer_orden}.
\end{proof}

Decimos que $\bar{x} \in \Rn$ es un \textbf{punto estacionario} del problema \eqref{eq:problema_general_restricciones} si $\bar{x} \in D$ y vale \eqref{eq:condicion_primal_primer_orden}. 

\begin{corollary}
    Sea $\bar{x} \in \interior D$ una solución local del problema \eqref{eq:problema_general_restricciones} tal que $f$ es diferenciable en $\bar{x}$. Entonces $\nabla f(\bar{x}) = 0$.
\end{corollary}

La condición de optimalidad puede transformarse a la forma primal-dual usando el cono dual.

\begin{definition}[Cono dual]\label{def:cono_dual}
    El \textbf{cono dual} de un cono $K \subset \Rn$ se define por
    \[
        K^* = \{y \in \Rn \mid \interno{y}{d} \le 0 \quad \forall d \in K\}.
    \]
\end{definition}

Usando esta noción, la condición de optimalidad \eqref{eq:condicion_primal_primer_orden} es equivalente a
\begin{equation}
    -\nabla f(\bar{x}) \in (\mathcal{B}_D(\bar{x}))^*. \label{eq:condicion_dual_primal}
\end{equation}

\begin{center}
\begin{tikzpicture}[>=Stealth, scale=0.9]
    \fill[MidnightBlue!10, draw=MidnightBlue!60!black, thick] (0,0) -- (2,-1) -- (2,1) -- cycle;
    \node[MidnightBlue!80!black] at (1.5,0.2) {$\mathcal{B}_D(\bar{x})$};
    \fill[red!10, draw=red!60!black, thick] (0,0) -- (-1,2) -- (1,2) -- cycle;
    \node[red!60!black] at (0,1.5) {$(\mathcal{B}_D(\bar{x}))^*$};
    \draw[->, thick, ForestGreen] (0,0) -- (0.8,0.2) node[right] {$\nabla f(\bar{x})$};
    \draw[->, thick, purple!80!black] (0,0) -- (-0.8,1.0) node[left] {$-\nabla f(\bar{x})$};
    \draw[dashed, gray] (0,0) -- (2,2);
    \draw[dashed, gray] (0,0) -- (-2,2);
\end{tikzpicture}
\end{center}

\begin{theorem}[Condición suficiente en forma primal]\label{thm:cs_primal}
    Sean $D \subset \Rn$ y $f : \Rn \to \R$ diferenciable en $\bar{x} \in D$. Entonces
    \begin{equation}
        \interno{\nabla f(\bar{x})}{d} > 0 \quad \forall d \in \mathcal{B}_D(\bar{x}) \setminus \{0\} \label{eq:cs_primal}
    \end{equation}
    si, y solamente si, $f$ satisface en $D$ en torno a $\bar{x}$ la condición de crecimiento lineal:
    \begin{equation}
        f(x) - f(\bar{x}) \ge \beta \norm{x - \bar{x}} \quad \forall x \in D \cap U. \label{eq:condicion_crecimiento_lineal}
    \end{equation}
    En particular, \eqref{eq:condicion_crecimiento_lineal} implica que $\bar{x}$ es minimizador local estricto.
\end{theorem}

\begin{proof}
    Supongamos que \eqref{eq:cs_primal} valga, pero no \eqref{eq:condicion_crecimiento_lineal}. Existe $\{x^k\} \subset D \setminus \{\bar{x}\}$ convergiendo a $\bar{x}$ tal que $f(x^k) - f(\bar{x}) < \frac{1}{k} \norm{x^k - \bar{x}}$. Definiendo $t_k = \norm{x^k - \bar{x}}$ y $d^k = (x^k - \bar{x})/t_k$, podemos extraer una subsucesión convergente $d^k \to d$ ($\norm{d}=1$). Así $d \in \mathcal{B}_D(\bar{x}) \setminus \{0\}$. Evaluando el gradiente, obtenemos $0 \ge \interno{\nabla f(\bar{x})}{d}$, en contradicción con \eqref{eq:cs_primal}.
    
    El recíproco es directo evaluando \eqref{eq:condicion_crecimiento_lineal} a lo largo de las trayectorias que definen el cono de Bouligand y tomando el límite para el gradiente.
\end{proof}

\begin{proposition}[Relaciones entre conos]\label{prop:relaciones_conos}
    Sean $D \subset \Rn$ y $\bar{x} \in D$. Los conos $\mathcal{T}_D(\bar{x})$ y $\mathcal{B}_D(\bar{x})$ son cerrados y se tiene que
    \[
        \{0\} \subset \cl \mathcal{V}_D(\bar{x}) \subset \mathcal{T}_D(\bar{x}) \subset \mathcal{B}_D(\bar{x}).
    \]
\end{proposition}